% Reviewed translation: PDF pages 361--365 / printed pages 347--351,
% plus the conclusion of Section 5 on PDF 366 / printed page 352.
\MathBookSourcePage{361}
此外, 当
$\phi\in H^{l+3/2}(\Omega)\cap W^{l+1,\infty}(\Omega)$ 时,
Lemma~\ref{lem:incompressible-further-improved-kernel-approximation} 给出最佳估计
\[
\|\pi_hv-v\|_X
\leq Ch^{l-1/2}
(\|\phi\|_{l+3/2,\Omega}+\|\phi\|_{l+1,\infty,\Omega}).
\]

现在转向 \eqref{eq:navier-branch-gh-consistency}. 有
\[
\langle G_h(\pi_hu)-G(u),v_h\rangle
=\widetilde a^{\pi_hu}(\pi_hu;\pi_hu,\boldsymbol v_h)
-a_1(u;u,\boldsymbol v_h)
\quad\forall\boldsymbol v_h=\operatorname{curl}\phi_h,
\quad\phi_h\in\Phi_h.
\]
简单重排各项得到
\[
\langle G_h(\pi_hu)-G(u),v_h\rangle
=\widetilde a^{\pi_hu}(\pi_hu;\pi_hu-u,\boldsymbol v_h)
+a_1(\pi_hu-u;u,\boldsymbol v_h).
\]
因此 Lemmas~\ref{lem:navier-upwind-consistency-bound},
\ref{lem:navier-upwind-form-bounded} 和
\eqref{eq:navier-upwind-mesh-dependent-norm} 一起给出
\begin{equation}\label{eq:navier-upwind-gh-consistency-bound}
\|G_h(\pi_hu)-G(u)\|_h
\leq Ch^{1/2}
(\|\psi\|_{2,\Omega}+|\mathring P_h\psi|_{1,4,\Omega})
[\mathring P_h\psi-\psi]_4.
\end{equation}
所以 \eqref{eq:navier-branch-gh-consistency} 由
Lemma~\ref{lem:navier-upwind-projection-mesh-norm} 得出. 此外, 由该 Lemma 和
\eqref{eq:navier-upwind-gh-consistency-bound} 推得
\begin{equation}\label{eq:navier-upwind-gh-consistency-rate}
\|G_h(\pi_hu)-G(u)\|_h
\leq Ch^{k-1/2}|\psi|_{2,\Omega}\|\psi\|_{k+1,\Omega}
\quad\forall u=(\operatorname{curl}\psi,\omega)\in V,
\quad\psi\in H^{k+1}(\Omega),\quad1\leq k\leq l.
\end{equation}

关于 \eqref{eq:navier-branch-generalized-gradient-consistency}, 应用
Definition~\ref{def:navier-upwind-generalized-gradient} 和
\eqref{eq:navier-upwind-regularized-derivative}:
\[
\begin{aligned}
\langle(\nabla G_h(\pi_hu)-DG(u))\cdot v_h,w_h\rangle
={}&\widetilde a^{\pi_hu}(\pi_hu;v_h,\boldsymbol w_h)
+\widetilde a^{\pi_hu}(v_h;\pi_hu,\boldsymbol w_h)\\
&-a_1(u;Pv_h,\boldsymbol w_h)-a_1(v_h;u,\boldsymbol w_h)
\quad\forall u\in V\\
={}&\widetilde a^{\pi_hu}(\pi_hu;v_h-Pv_h,\boldsymbol w_h)
+\widetilde a^{\pi_hu}(v_h;\pi_hu-u,\boldsymbol w_h)\\
&+a_1(\pi_hu-u;Pv_h,\boldsymbol w_h),
\end{aligned}
\]
这里考虑了 $u$ 和 $Pv_h$ 的正则性. 因此由 Lemmas~
\ref{lem:navier-upwind-boundary-almost-lipschitz} 和
\ref{lem:navier-upwind-consistency-bound} 得到
\[
\begin{aligned}
\| (\nabla G_h(\pi_hu)-DG(u))\cdot v_h\|_h
\leq C\{h^{1/2}(&|\mathring P_h\psi|_{1,4,\Omega}[\phi_h-\phi]_4\\
&+|v_h|[\mathring P_h\psi-\psi]_4)
+|v_h||\mathring P_h\psi-\psi|_{1,4,\Omega}\},
\end{aligned}
\]
其中 $Pv_h=(\operatorname{curl}\phi,\theta_h)$,
$u=(\operatorname{curl}\psi,\omega)\in V$, 且
$v_h=(\operatorname{curl}\phi_h,\theta_h)\in V_h$. 因而
Lemma~\ref{lem:navier-upwind-projection-mesh-norm} 给出
\[
\|(\nabla G_h(\pi_hu)-DG(u))\cdot v_h\|_h
\leq Ch^{1/2}(|\psi|_{1,4,\Omega}+\|\psi\|_{2,\Omega})|v_h|,
\]
这蕴含 \eqref{eq:navier-branch-generalized-gradient-consistency}.

最后还需验证 \eqref{eq:navier-branch-gh-generalized-lipschitz}.
取 $u_h,u_h^*$ 和 $u_h^0\in V_h$; 根据定义, 对所有
$\boldsymbol w_h=\operatorname{curl}\chi_h$, $\chi_h\in\Phi_h$, 有
\[
\begin{aligned}
&\langle G_h(u_h)-G_h(u_h^*)
-\nabla G_h(u_h^0)\cdot(u_h-u_h^*),w_h\rangle\\
={}&\widetilde a^{u_h}(u_h;u_h,\boldsymbol w_h)
-\widetilde a^{u_h^*}(u_h^*;u_h^*,\boldsymbol w_h)\\
&-\widetilde a^{u_h^0}(u_h^0;u_h-u_h^*,\boldsymbol w_h)
-\widetilde a^{u_h^0}(u_h-u_h^*;u_h^0,\boldsymbol w_h)\\
={}&\widetilde a^{u_h^0}(u_h-u_h^0;u_h-u_h^*,\boldsymbol w_h)
+\widetilde a^{u_h^0}(u_h-u_h^*;u_h^*-u_h^0,\boldsymbol w_h)\\
&+\{a_2^{u_h}(u_h;u_h,\boldsymbol w_h)
-a_2^{u_h^0}(u_h;u_h,\boldsymbol w_h)
-a_2^{u_h^*}(u_h^*;u_h^*,\boldsymbol w_h)
+a_2^{u_h^0}(u_h^*;u_h^*,\boldsymbol w_h)\}.
\end{aligned}
\]

\MathBookSourcePage{362}
大括号中的表达式可重排为
\[
\begin{aligned}
&\{a_2^{u_h}(u_h-u_h^*;u_h,\boldsymbol w_h)
-a_2^{u_h^0}(u_h-u_h^*;u_h,\boldsymbol w_h)\}\\
&+\{a_2^{u_h^*}(u_h^*;u_h-u_h^*,\boldsymbol w_h)
-a_2^{u_h^0}(u_h^*;u_h-u_h^*,\boldsymbol w_h)\}\\
&+\{a_2^{u_h}(u_h^*;u_h,\boldsymbol w_h)
-a_2^{u_h^*}(u_h^*;u_h,\boldsymbol w_h)\}.
\end{aligned}
\]
反复使用 Lemma~\ref{lem:navier-upwind-boundary-almost-lipschitz}, 可将该表达式控制为
\[
\begin{aligned}
C_1h^{1/2}\{&|u_h||\psi_h-\psi_h^*|_{1,4,\Omega}
+|\psi_h^*-\psi_h^0|_{1,4,\Omega}|u_h-u_h^*|\}|w_h|\\
\leq{}&C_2h^{1/2}|u_h-u_h^*|
\{|u_h-u_h^0|+|u_h^*-u_h^0|+|u_h^0|\}|w_h|
\quad\forall w_h\in V_h.
\end{aligned}
\]
因此 Lemma~\ref{lem:navier-upwind-form-bounded} 给出界
\[
\begin{aligned}
&\|G_h(u_h)-G_h(u_h^*)
-\nabla G_h(u_h^0)\cdot(u_h-u_h^*)\|_h\\
&\qquad\leq C_3|u_h-u_h^*|
\{|u_h-u_h^0|+|u_h^*-u_h^0|+h^{1/2}|u_h^0|\}.
\end{aligned}
\]
这以函数
\[
L_h(\mu;v)=C_4(\mu+h^{1/2}v)
\]
蕴含 \eqref{eq:navier-branch-gh-generalized-lipschitz}.
显然, $L_h$ 连续且关于每个变量单调递增, 并且
\[
\lim_{h\to0}L_h(0;v)=\lim_{h\to0}(h^{1/2}v)=0
\quad\forall v\in\mathbb R_+.
\]

Theorem~\ref{thm:navier-branch-nondifferentiable-approximation} 的全部假设都已满足,
所以可把其结论应用于上风格式 \eqref{eq:navier-upwind-discrete-problem},
从而得到其解分支的存在性, 唯一性和收敛性. 此外, 比较误差估计
\eqref{eq:navier-branch-nondifferentiable-error} 和
\eqref{eq:navier-branch-fundamental-error}, 并考虑
\eqref{eq:navier-upwind-gh-consistency-rate}, 容易推知上风格式的误差界与中心格式完全相同.
换言之, 在这种情形下, \textit{上风化不改变格式的阶}. 下面的 Theorem 总结这些结果.

\setcounter{statement}{0}
\begin{Theorem}\label{thm:navier-upwind-branch-error}
设 $\Omega$ 是有界凸多边形, 且 $\mathcal T_h$ 是 $\bar\Omega$ 的一致正则三角剖分族.
对 $\boldsymbol f\in L^t(\Omega)^2$, $t\in[r,2)$, 设
\[
\{(\lambda,(\operatorname{curl}\psi(\lambda),\omega(\lambda)));
\ \lambda=1/\nu\in\Lambda\}
\]
是 Navier--Stokes Problem~\eqref{eq:navier-upwind-continuous-problem}
的一个非奇异解分支. 则对充分小的 $h\leq h_0$, 存在唯一的上风格式
\eqref{eq:navier-upwind-discrete-problem} 的 $C^0$ 解分支
\[
\{(\lambda,(\operatorname{curl}\psi_h(\lambda),\omega_h(\lambda)));
\ \lambda=1/\nu\in\Lambda\},
\]
且满足误差估计
\begin{equation}\label{eq:navier-upwind-primary-error}
\sup_{\lambda\in\Lambda}\left\{
|\psi_h(\lambda)-\psi(\lambda)|_{1,s,\Omega}
+\|\omega_h(\lambda)-\omega(\lambda)\|_{0,\Omega}
\right\}
\leq C_1h^\alpha,
\qquad
\alpha=\begin{cases}
1/\gamma,&l=1,\\
2/\gamma,&l\geq2,
\end{cases}
\quad \frac1t+\frac1\gamma=1,
\end{equation}
其中常数 $C$ 与 $h$ 和 $\lambda$ 无关.
\end{Theorem}

\MathBookSourcePage{363}
当 $t=2$ 且 $l\geq2$ 或 $\psi\in W^{2,\infty}(\Omega)$ 时, 这个界仍然成立.
若 $l=1$ 且 $t=2$, 则 \eqref{eq:navier-upwind-primary-error} 的左端由
\[
C_2(\varepsilon)h^{1/2-\varepsilon}
\quad\forall\varepsilon>0
\]
控制.

此外, 若对某个实数 $m\in[1,l-1/2]$, 映射
\[
\lambda\mapsto\psi(\lambda)
\]
从 $\Lambda$ 连续地映到
$H^{m+2}(\Omega)\cap W^{m+3/2,\infty}(\Omega)$, 则
\begin{equation}\label{eq:navier-upwind-higher-regularity-error}
\sup_{\lambda\in\Lambda}\left\{
|\psi_h(\lambda)-\psi(\lambda)|_{1,s,\Omega}
+\|\omega_h(\lambda)-\omega(\lambda)\|_{0,\Omega}
\right\}
\leq C_3h^m.
\end{equation}

最后, 与 Subsection~\ref{subsec:navier-centered-stream-vorticity} 极为相似的论证
可以细化 $|\psi-\psi_h|_{1,\Omega}$ 的误差估计, 并得到与
Theorem~\ref{thm:navier-centered-stream-h1-error} 相同的收敛阶.

\setcounter{statement}{1}
\begin{Theorem}\label{thm:navier-upwind-stream-h1-error}
设 $\Omega$ 和 $\mathcal T_h$ 如 Theorem~\ref{thm:navier-upwind-branch-error},
并假设 Navier--Stokes Problem~\eqref{eq:navier-upwind-continuous-problem}
的非奇异解分支具有正则性
\[
\lambda\mapsto\psi(\lambda)\in
\begin{cases}
C^0(\Lambda;H^{l+1}(\Omega)),&l\geq2,\\
C^0(\Lambda;H^3(\Omega)),&l=1.
\end{cases}
\]
则上风格式 \eqref{eq:navier-upwind-discrete-problem} 的逼近解 $\psi_h$ 满足
\begin{equation}\label{eq:navier-upwind-stream-h1-optimal}
|\psi(\lambda)-\psi_h(\lambda)|_{1,\Omega}
\leq
\begin{cases}
C_1h^l,&l\geq2,\\
C_2(\varepsilon)h^{1-\varepsilon},&l=1,\quad\forall\varepsilon>0,
\end{cases}
\end{equation}
其中常数与 $h$ 和 $\lambda$ 无关.
\end{Theorem}
\begin{Proof}
任取 $\boldsymbol g\in L^2(\Omega)^2$, 并引入与
\eqref{eq:navier-centered-stream-dual-linearized} 类似的对偶线性化 Navier--Stokes 问题:
求 $z=(\boldsymbol z=\operatorname{curl}\chi,\eta)\in V$, 使
\begin{equation}\label{eq:navier-upwind-dual-linearized}
(v,\theta)+\lambda a_1(v;u(\lambda),\boldsymbol z)
-\lambda a_1(u(\lambda);z,\boldsymbol v)
=(\boldsymbol g,\operatorname{curl}\phi)
\quad\forall v=(\boldsymbol v=\operatorname{curl}\phi,\theta)\in V.
\end{equation}
一方面已知
\[
a_1(u;z,\boldsymbol v)=-a_1(u;v,\boldsymbol z),
\]
另一方面, 对所有 $u,v,z\in V$,
\[
a_1(u;v,\boldsymbol z)+a_1(v;u,\boldsymbol z)
=\langle DG(u)\cdot v,z\rangle.
\]
因此 Problem~\eqref{eq:navier-upwind-dual-linearized} 与
Problem~\eqref{eq:navier-centered-stream-dual-linearized} 相同,
故其解具有 Lemma~\ref{lem:navier-centered-stream-dual-regularity} 中陈述的正则性.

现在像 Theorem~\ref{thm:navier-centered-stream-h1-error} 那样论证, 容易得到
\[
\begin{aligned}
(\boldsymbol g,\operatorname{curl}(\psi-\psi_h))
={}&(\eta-\eta_h,\omega-\omega_h)
+\widetilde b(u-u_h,v-\theta_h)
+\widetilde b(z-z_h,\omega-\mu_h)\\
&+\lambda a_1(u-u_h;u,\boldsymbol z)
-\lambda a_1(u;z,u-u_h)
-\lambda a_1(u;u,\boldsymbol z_h)\\
&+\lambda a_1(u_h;u_h,\boldsymbol z_h)
+\lambda a_2^{u_h}(u_h;u_h,\boldsymbol z_h)
\end{aligned}
\]
对所有 $z_h=(\boldsymbol z_h=\operatorname{curl}\chi_h,\eta_h)\in V_h$
和所有 $\mu_h,\theta_h\in\Theta_h$ 成立.

\MathBookSourcePage{364}
因此, 对各项作简单整理并利用 \eqref{eq:navier-upwind-integration-by-parts}, 得到
\begin{equation}\label{eq:navier-upwind-dual-error-identity}
\begin{aligned}
(\boldsymbol g,\operatorname{curl}(\psi-\psi_h))
={}&(\eta-\eta_h,\omega-\omega_h)
+\widetilde b(u-u_h,v-\theta_h)
+\widetilde b(z-z_h,\omega-\mu_h)\\
&+\lambda\{a_1(u-u_h;u,\boldsymbol z-\boldsymbol z_h)
+\widetilde a^{u_h}(u_h;u-u_h,\boldsymbol z-\boldsymbol z_h)\\
&\hspace{28mm}-a_1(u-u_h;z,u-u_h)\}.
\end{aligned}
\end{equation}
但该证明中精细的一步, 是适当估计乘在 $\lambda$ 前的因子中的
$\widetilde a(\mathord\cdot;\mathord\cdot,\mathord\cdot)$ 项. 若应用
Lemma~\ref{lem:navier-upwind-form-bounded}, 得到
\[
|\widetilde a^{u_h}(u_h;u-u_h,\boldsymbol z-\boldsymbol z_h)|
\leq Ch^{1/2}|\psi_h|_{1,4,\Omega}
[\psi-\psi_h]_2[\chi-\chi_h]_2.
\]
当 $l=1$ 时, 这个上界没有用, 因为
$\inf_{\chi_h}[\chi-\chi_h]_2=O(1)$. 取而代之, 最好利用
$\chi\in H^3(\Omega)$ 的事实, 把前一上界替换为
\[
\begin{aligned}
|\widetilde a^{u_h}(u_h;u-u_h,\boldsymbol z-\boldsymbol z_h)|
\leq C_1|\psi_h|_{1,4,\Omega}[\psi-\psi_h]_2
\Bigg\{&|\chi-\chi_h|_{1,4,\Omega}\\
&+h\left(\sum_{\kappa\in\mathcal T_h}
|\chi-\chi_h|_{2,4,\kappa}^4\right)^{1/4}\Bigg\}.
\end{aligned}
\]
因此, 若选择由 \eqref{eq:navier-upwind-projection-operator} 定义的
$z_h=\pi_hz\in V_h$, 则 $\chi_h=\mathring P_h\chi$ 蕴含
\[
|\widetilde a^{u_h}(u_h;u-u_h,\boldsymbol z-\boldsymbol z_h)|
\leq C_2h|\psi_h|_{1,4,\Omega}[\psi-\psi_h]_2
\|\chi\|_{2,4,\Omega}.
\]
最后注意
\begin{equation}\label{eq:navier-upwind-stream-mesh-error}
\begin{aligned}
[\psi-\psi_h]_2
&\leq[\psi-\mathring P_h\psi]_2
+[\mathring P_h\psi-\psi_h]_2\\
&\leq C_3h^{l-1}\|\psi\|_{l+1,\Omega}
+C_4|\pi_hu-u_h|,
\end{aligned}
\end{equation}
这里使用了 \eqref{eq:navier-upwind-projection-mesh-estimate},
\eqref{eq:navier-upwind-uniform-broken-h2} 和
\eqref{eq:navier-upwind-projection-operator}. 这给出形如
\[
|\widetilde a^{u_h}(u_h;u-u_h,\boldsymbol z-\boldsymbol z_h)|
\leq C_5\|\psi\|_{2,\Omega}\|\chi\|_{2,4,\Omega}
(h^l\|\psi\|_{l+1,\Omega}+h|\pi_hu-u_h|)
\]
的上界. \eqref{eq:navier-upwind-dual-error-identity} 中其他项容易估计,
证明完全像 Theorem~\ref{thm:navier-centered-stream-h1-error} 那样结束.
\end{Proof}

\subsection{用上风格式逼近压力}
\label{subsec:navier-upwind-pressure}

在 Subsection~\ref{subsec:navier-centered-stream-vorticity} 中已经看到,
Navier--Stokes 系统 \eqref{eq:navier-upwind-continuous-problem} 所对应的压力项 $p$
是下列问题的解: 求 $p\in W^{1,r}(\Omega)\cap L_0^2(\Omega)$, 使
\begin{equation}\label{eq:navier-upwind-continuous-pressure}
(\operatorname{grad}p,\operatorname{grad}q)
=\left(\boldsymbol f-\nu\operatorname{curl}\omega
-\sum_{j=1}^2u_j\partial\boldsymbol u/\partial x_j,
\operatorname{grad}q\right)
\quad\forall q\in W^{1,s}(\Omega).
\end{equation}
类似地, 为恢复与上风格式 \eqref{eq:navier-upwind-discrete-problem} 相联系的压力 $p_h$,
引入由 \eqref{eq:navier-centered-stream-discrete-spaces} 定义的空间 $Q_h$.

\MathBookSourcePage{365}
即
\begin{equation}\label{eq:navier-upwind-pressure-space}
Q_h=\{q_h\in\mathcal C^0(\bar\Omega)\cap L_0^2(\Omega);
\ q_h|_\kappa\in P_k\ \forall\kappa\in\mathcal T_h\},
\qquad k=\min(1,l-1),
\end{equation}
并把 \eqref{eq:navier-upwind-continuous-pressure} 离散为: 求 $p_h\in Q_h$, 使
\begin{equation}\label{eq:navier-upwind-discrete-pressure}
(\operatorname{grad}p_h,\operatorname{grad}q_h)
=(\boldsymbol f-\nu\operatorname{curl}\omega_h,
\operatorname{grad}q_h)
-\widetilde a^{u_h}(u_h;u_h,\operatorname{grad}q_h)
\quad\forall q_h\in Q_h.
\end{equation}
显然, 这个问题有唯一解.

为估计误差 $p-p_h$, 使用与 Theorem~\ref{thm:incompressible-sfvp-pressure-error}
相同的对偶性论证. 引入由下式定义的函数
$\boldsymbol v\in H^1(\Omega)^2$:
\[
\operatorname{div}\boldsymbol v=p-p_h,
\qquad
|\boldsymbol v|_{1,\Omega}\leq C_1\|p-p_h\|_{0,\Omega},
\]
并把它分解为
\[
\boldsymbol v=\operatorname{grad}q+\operatorname{curl}\phi.
\]
由于假设 $\Omega$ 是凸多边形, $q$ 和 $\phi$ 都属于 $H^2(\Omega)$, 且
\[
\|q\|_{2,\Omega}+\|\phi\|_{2,\Omega}
\leq C_2|\boldsymbol v|_{1,\Omega}
\quad\text{(参见 Lemma~\ref{lem:incompressible-sfvp-hypothesis-h2})}.
\]
此外, 若三角剖分 $\mathcal T_h$ 一致正则,
Lemma~\ref{lem:incompressible-sfvp-hypothesis-h2} 表明
\begin{equation}\label{eq:navier-upwind-pressure-projection-error}
|q-P_hq|_{1,s,\Omega}+|\phi-I_h\phi|_{1,s,\Omega}
\leq C_3h^{2/s}|\boldsymbol v|_{1,\Omega}.
\end{equation}
现在像 Theorem~\ref{thm:incompressible-sfvp-pressure-error} 中那样可写成
\[
\begin{aligned}
\|p-p_h\|_{0,\Omega}^2
={}&(\operatorname{grad}(q_h-p),\operatorname{grad}(q-P_hq))\\
&+(\operatorname{grad}(p_h-p),\operatorname{grad}(P_hq))
\qquad\forall q_h\in Q_h.
\end{aligned}
\]
因此, 为估计 $p-p_h$, 必须对第二项导出一个精细界.
从 \eqref{eq:navier-upwind-discrete-pressure} 中减去
\eqref{eq:navier-upwind-continuous-pressure}, 得到
\[
\begin{aligned}
&(\operatorname{grad}(p_h-p),\operatorname{grad}(P_hq))\\
={}&-\nu(\operatorname{curl}(\omega_h-\omega),\operatorname{grad}(P_hq))
+a_1(u;u,\operatorname{grad}(P_hq))
-\widetilde a^{u_h}(u_h;u_h,\operatorname{grad}(P_hq))\\
={}&\nu(\operatorname{curl}(\omega-\omega_h),\boldsymbol v_h-\boldsymbol v)
+\nu(\omega-\omega_h,\operatorname{curl}\boldsymbol v)\\
&+\widetilde a^{u_h}(u_h;u-u_h,\boldsymbol v_h-\boldsymbol v)
+a_1(u-u_h;u,\boldsymbol v_h-\boldsymbol v)\\
&+a_1(u_h-u;\boldsymbol v,u_h)+a_1(u;\boldsymbol v,u_h-u),
\end{aligned}
\]
其中
\[
\boldsymbol v_h=\operatorname{grad}(P_hq)+\operatorname{curl}(I_h\phi).
\]
因此 \eqref{eq:navier-upwind-pressure-projection-error} 与对
$\widetilde a^{u_h}(\mathord\cdot;\mathord\cdot,\mathord\cdot)$ 和
$a_1(\mathord\cdot;\mathord\cdot,\mathord\cdot)$ 的熟悉估计一起,
给出与 \eqref{eq:incompressible-sfvp-pressure-error} 类似的结果.

\setcounter{statement}{4}
\begin{Lemma}\label{lem:navier-upwind-pressure-error}
设 $\Omega$ 是有界凸多边形, 且 $\mathcal T_h$ 是 $\bar\Omega$ 的一致正则三角剖分.
若对某个实数 $t\in[r,2]$, $p$ 和 $\omega$ 属于 $W^{1,t}(\Omega)$,
则 $p$ 的误差为

\MathBookSourcePage{366}
\begin{equation}\label{eq:navier-upwind-pressure-error}
\begin{aligned}
\|p-p_h\|_{0,\Omega}\leq C\Bigg\{
&h^{2/\gamma}\inf_{q_h\in Q_h}|p-q_h|_{1,t,\Omega}
+\nu\|\omega-\omega_h\|_{0,\Omega}\\
&+\nu\inf_{\theta_h\in\Theta_h}
(\|\omega-\theta_h\|_{0,\Omega}
+h^{2/\gamma}|\omega-\theta_h|_{1,t,\Omega})\\
&+h^{1/2}[\psi-\psi_h]_4
(|\psi_h|_{1,4,\Omega}+|\psi|_{2,\Omega})
\Bigg\},
\end{aligned}
\end{equation}
其中 $1/\gamma+1/t=1$,
$u_h=(\boldsymbol u_h=\operatorname{curl}\psi_h,\omega_h)$ 和
$u=(\boldsymbol u=\operatorname{curl}\psi,\omega)$ 分别表示
\eqref{eq:navier-upwind-discrete-problem} 和
\eqref{eq:navier-upwind-continuous-problem} 的解.
\end{Lemma}

由 \eqref{eq:navier-upwind-pressure-error} 和
\eqref{eq:navier-upwind-stream-mesh-error} 容易看出, 压力 $p_h$ 与速度具有相同收敛阶.

\setcounter{statement}{2}
\begin{Theorem}\label{thm:navier-upwind-pressure-convergence}
在 Theorem~\ref{thm:navier-upwind-branch-error} 的假设下,
由 Problem~\eqref{eq:navier-upwind-discrete-pressure} 定义的压力 $p_h$
以与速度相同的精度阶收敛到 $p$. 即若
$\boldsymbol f\in L^t(\Omega)^2$, $t\in[r,2)$, 则
\[
\sup_{\lambda\in\Lambda}\|p(\lambda)-p_h(\lambda)\|_{0,\Omega}
\leq C_1h^\alpha,
\qquad
\alpha=\begin{cases}
1/\gamma,&l=1,\\
2/\gamma,&l=2,
\end{cases}
\]
其中 $1/\gamma+1/t=1$. 当 $\psi\in W^{2,\infty}(\Omega)$ 或 $l\geq2$ 时,
该界可以推广到 $t=2$. 当 $l=1$ 且 $t=2$ 时,
\[
\sup_{\lambda\in\Lambda}\|p(\lambda)-p_h(\lambda)\|_{0,\Omega}
\leq C_2(\varepsilon)h^{1-\varepsilon}.
\]
最后, 若对某个实数 $m\in[1,l-1/2]$, 映射
\[
\lambda\mapsto(\psi(\lambda),p(\lambda))
\]
从 $\Lambda$ 连续地映到
\[
[H^{m+2}(\Omega)\cap W^{m+3/2,\infty}(\Omega)]\times H^m(\Omega),
\]
则
\[
\sup_{\lambda\in\Lambda}\|p(\lambda)-p_h(\lambda)\|_{0,\Omega}
\leq C_3h^m.
\]
\end{Theorem}
